\scanpage{79}
\mspage{72}
% archive sheet number: 78
Le dernier énoncé du $2^\circ$ du th.-scholie résulte du calcul de $u_{\unsure{*|\circ}}^W$, compte tenu que $\deleted{\unsure{i'|I'}\longleftarrow\unsure{s|*}}\ \deleted{\unsure{s\backslash}}I'_{/s}\unsure{-|\longrightarrow}I'$ est un iso local, donc le \inserted{foncteur} image inverse \deleted{\unsure{d'un|d'une} \unsure{est|et}} envoie
\[
M(I')_f\unsure{-|\longrightarrow}M(I'_{/s})_f,
\]
donc si $\unsure{G'|C'}$ est dans $M(I')_f$,
\[
u_*(G')(s)\simeq p_{s*}(G'|I'_{/s})
\]
est un calcul de $p_{s*}^W(G'|I'_{/s})$.

Cela achève la démonstration du th.-scholie sur les images directes.

\uline{Commentaire.} 1) Ce théorème me paraît loin d'être satisfaisant, c'est un pis-aller. J'avais espéré, avec les hypothèses du \S\ 3, qu'on aurait l'existence du $u_*^W$ pour \uline{toute} flèche $u:I'\to I$ dans Cat. Je doute à présent que ce soit le cas, mais n'ai pas encore \unsure{trouvé|essayé} de faire un contre-exemple. Les cas les plus importants \inserted{pour la question d'admissibilité d'une catégorie $I$} me paraissent les suivants:

a) Catégorie \inserted{$I$} discrète infinie (disons dénombrable). La question revient ici à celle-ci: soient $f_i:X_i\to Y_i$ dans $W$, avec $X_i,Y_i$ fibrants donc dans $F$. Alors $f=\prod f_i:\prod X_i\to\prod Y_i$ est-il dans $W$?
